% Section 13: discussion, limitations, conclusion (P4)
\subsection{Where the stack itself fails}
The honest reading of our own frame is that cross-layer engineering buys no reliability without
instrumentation, and that two failure modes recur across the case set and bound what the frame can
promise. First, per-dimension
certification does not compose: a validated contract plus a calibrated gate plus a grounded
retriever can still fail on their \emph{interaction}, and no anchor in our evidence base
measures that composition (cases 1 and 3; \cref{sec:case-rag}, \cref{sec:case-hybrid}).
Second, verification mechanisms are themselves attackable or miscalibrated
\cite{id260601066,id260930217} (cases 6 and 5), so the assurance layer needs a threat model of
its own.

\subsection{What the field still lacks}
Eight gaps remain open, and each is phrased as a testable ask:
(i) \emph{end-to-end system evaluations} of assembled stacks (harness + gate + escalation)
under adversarial conditions;
(ii) \emph{certified composition} results for chained safeguards (conformal $\circ$ validators
$\circ$ grounding), which extend the [T]-tier machinery to mismatches between layers;
(iii) \emph{propagation benchmarks} for hybrid specialist stacks, where the Case 3 gap
(\cref{sec:case-hybrid}) is the cleanest measurement target our evidence set exposes;
(iv) \emph{language-native latent-prediction training signals} that survive the
conditional-concentration objection (JEPA paradox), with assumptions actually testable;
(v) \emph{decision-model interface stability}: order robustness and option-set scaling are
measurable, and current evidence suggests they are not incidental \cite{id261002267};
(vi) \emph{reconstruction-fidelity measurement for staged-tokenizer generators}: a test that
separates token-exactness from meaning-equivalence, so a contract can state which one it
enforces;
(vii) \emph{escalation-queue measurement}: precision, recall and minutes per review for
reviewers working a production LLM escalation queue, which no study in our sweep reports
(\cref{sec:l5-calibrated});
(viii) \emph{a configuration record for deployment}: the engine, precision and co-tenancy a
score was measured under, published with the cutoff fitted on it (\cref{sec:l5-bridge}).

\subsection{What we refuse to compare}
Layer identities are inherited envelopes, not benchmark scores; we decline direct comparisons
across families (e.g., discriminative vs.\ generative ``hallucination rates'') when the error
currencies are not measurable on a shared task. Where the evidence base has no shared-instrument
anchor we say so, and where it has one, each row names the instrument it rests on rather than pooling
across instruments (\cref{tab:crosscase}).

\subsection{Limitations}
Scope is LLM-centric (text and typed-decision workloads); multimodal and embodied control appear
only as instances of an architecture. The paper is survey-grade: every claim
either resolves to a public artifact we checked or is tagged where checking stops --- [W] pending,
[N]/[GAP] no anchor --- and [V]/[A] mark evidence that is documented but not independently
reproduced. We run no new benchmarks; contested rows in
\cref{fig:evidence} remain contested at submission time and are labeled as such. The seven case
studies cover the stack's decision space but not exhaustively; the exclusion register applies
to our choices as it did to the candidates'.

The three validity bounds --- instrument, contamination, precision --- are fixed in
\cref{sec:protocol}, and they apply to our own anchors as much as to the sources we cite. What is
ours to disclose is narrower: our empirical numbers are point estimates from single harnesses, and
a point estimate is not a difference \cite{errorbars2024,evalci2026}.

The boundary to standards and regulation is a division of labor, not a crosswalk. The NIST AI
RMF 1.0 \cite{nistairmf} (January 2023, voluntary, and under revision as of this writing) and
its Generative AI Profile (NIST AI 600-1 \cite{nistgai600}, July 2024) specify risk-governance
functions: govern, map, measure, manage. ISO/IEC 42001:2023 \cite{iso42001} defines a
certifiable AI management system. The EU AI Act (Regulation (EU) 2024/1689 \cite{euaiact2024})
specifies provider and deployer duties; the digital-omnibus Regulation (EU) 2026/1744
\cite{omnibus2026} defers the Chapter~III high-risk obligations for Annex~III systems to December
2027 (Annex~I systems to August 2028), while the post-market monitoring and incident reporting
duties (Articles 72--73) stay on the 2026 timetable. None of these instruments rules on which layer absorbs a failure;
their unit is the organization, not the stack instance. They prescribe process: who must
document, assess, monitor, and report. A certificate attests to the process, never to the
residual error. Our five axes ask a different question: which stack instance to select for a
reliability-constrained workload. A management-system certificate does not answer it. The
relation runs one way. Our evidence classes can describe what a conformity assessment rests
on, and we do not translate controls into layers or claim that any certificate implies a
reliability property. A reader who needs compliance should read the instrument; a reader who
needs a selection decision should read \cref{sec:selection} and then the instrument.

\subsection{Four live fronts}
Four fronts are live, and we do not forecast which matures. Our claim is narrower: each would
change a named layer's contract, and none would replace the routing thesis that reliability is
assembled across L0--L5 and is not won inside any single model. Coarse-to-fine text generation
\cite{kumbong2026ncm,ncparch2026} is a Layer~1 objective family with a Layer~4 consequence: its
level order is a schedule (\cref{fig:schedule-atlas}) and the frozen tokenizer behind it is the
acceptance surface; the authors' own \evaurep{97M$\to$725M check still shows scale helping inside
the family}. Whether that hierarchy is learned, as in H-Net's dynamic chunking \cite{hnet2025}, or
fixed by a tokenizer remains an open schedule question; the atlas's three lanes differ in
commitment count and revisability, and learned chunking would move both.

Verifier-driven test-time training \cite{id250519475} closes a loop the stack now draws open:
adaptation at inference is only as reliable as the verifier choosing its pseudo-labels, so a
Layer~3 mechanism imports a Layer~5 obligation.

Risk-controlled serving \cite{id260933887} turns the L0 budget into an operating point rather
than a fixed constant: the budget rule of \cref{sec:selection} made executable, where
same-quality-faster buys candidates and checks rather than raw throughput.

Production diffusion serving \cite{inception2026mercury3,llada2026two} is the fourth. A hosted
text endpoint refines blocks in parallel over a number of passes the vendor sets, exposes no
step count, seed or token score, and floors temperature at 0.5, so the contract question becomes
which endpoint reports a score at all; the open-weight branch keeps
the sampler configurable. Mercury~3 postdates our cutoff, so we read it as vendor documentation
[V], not corpus evidence.

\subsection{Conclusion}
We asked when scaling a general-purpose model is the right lever and when a changed learning
objective, representation, decision interface, or verification pipeline is a better route to
reliable AI (\cref{sec:intro}). The evidence supports a narrower answer: both levers are real,
but they buy different currencies, and no currency is guaranteed by any single layer. Objective
families fail through design choices (an envelope they inherit), and the scaling story
that dominates current practice sells fluency and coverage without the contract, calibration and
assurance machinery that makes those properties deployable. Our contribution is a cross-layer
taxonomy (L0--L5) that names the decision point for each reliability claim, and a selection
framework that routes by task rather than by model preference. The claim ledger separates
theoretical guarantees, independent measurement, author- and vendor-reported results, synthesis,
and what remains genuinely contested, so that the next reliability argument starts from an
evidence map rather than a benchmark leaderboard \cite{leaderillusion2025}. Seven instances, one
of them an honest null, showed the same shape: the reliability-bearing move sat above the layer
where the failure originated, and the newest one, diffusion serving, moved by procurement at L4
rather than by engineering above the sampler. Configuration-relative reliability changes what a
cutoff means: an error rate belongs to the engine, precision and co-tenancy it was measured
under, so a stack change is a model change and the record has to say so. The residual is a
quantity the operator measures and owns.
